\documentclass[11pt]{article}
\usepackage{finprobs}

\title{Cross-Sectional Aggregation and the Scale-Dependence\\ of Roughness in log S-fBM}
\author{finprobs Research Project}
\date{July 2026}

\begin{document}
\maketitle

\begin{abstract}
Empirically, estimated Hurst roughness of log-volatility differs sharply by aggregation level: single stocks are estimated near $H\approx0.01$, broad indices near $H\approx0.1$. We ask whether this is a derivable consequence of cross-sectional aggregation, in the spirit of a renormalization-group flow, or merely a fitted regularity. We show a first-principles derivation attempt from linear mixing of heterogeneous log-vol processes stalls at a specific, identifiable circularity; we then show that a recent paper closes exactly that gap with a rigorous, checkable threshold condition -- but only conditional on an unexplained empirical input. We rule out a ``central-limit-theorem-of-central-limit-theorems'' shortcut to that input, propose (without endorsing) a speculative microstructural alternative, and independently confirm the aggregation mechanism itself by simulation and by fresh real-market estimation. We conclude the aggregation \emph{consequence} is now well-supported; the aggregation \emph{cause} remains open.
\end{abstract}

\section{Background and motivation}

In the log S-fBM model of \citet{wu2022rough}, log-volatility is a stationary Gaussian process with three free parameters $(H,\lambda^2,T)$ and covariance
\begin{equation}
C_\omega(\tau) = \frac{\nu^2}{2}\left[T^{2H}-\tau^{2H}\right], \qquad \nu^2 = \frac{\lambda^2}{H(1-2H)}. \label{eq:cov}
\end{equation}
\citet{gatheral2018volatility} estimate $H\approx0.1$ for major equity indices from high-frequency realized volatility, a result now widely cited as ``volatility is rough.'' Subsequent panel studies of individual names \citep{bennedsen2022decoupling,bolko2023gmm} find single-stock $H$ estimates that are smaller and more heterogeneous, often near $H\approx0$--$0.02$. \citet{wu2022rough} themselves flag the resulting question -- why does aggregation raise the effective roughness exponent? -- as an explicit open problem for future work, not something they resolve.

We ask: is $H_{\text{index}} > H_{\text{single-name}}$ a \emph{derivable} consequence of the model's own structure under cross-sectional aggregation, analogous to a renormalization-group flow of the deformation parameter, or is it merely an empirical regularity awaiting an explanation?

\section{Related work}

\citet{zarhali2025nested} (subsequently retitled ``A Nested Factor Model for Equity Markets: Reconciling Multifractal Stock Returns and Rough Index Volatilities'') model single-name log-volatility as
\begin{equation}
\omega_i(t) = \beta_i F(t) + X_i(t), \label{eq:factor}
\end{equation}
with a common factor $F$ of Hurst $H_F\approx0.11$ and idiosyncratic components $X_i$ that are ``super-rough''/near-multifractal ($H_i\approx0$), the split itself taken as an empirical input from \citet{wu2022rough}. Classical aggregation theory for heterogeneous processes \citep{granger1980long,zaffaroni2004aggregation} shows that mixing many short-memory or heterogeneous-persistence processes can alter the memory/self-similarity properties of the aggregate -- the general mechanism family this result belongs to.

\begin{finding}
\citet{gatheral2018volatility}'s $H\approx0.1$ estimate is the same underlying SPX-derived empirical anchor used both here and in the companion paper on SPX/VIX calibration (this series, Paper~4). It should be read as one empirical data point appearing in two contexts, not as two independent corroborations of rough volatility.
\end{finding}

\section{A stalled first-principles derivation}

We attempted to derive \eqref{eq:factor}'s consequence directly from linear aggregation, without presupposing the smooth/rough split.

\begin{proposition}[Fixed point under linear aggregation]
Let $\{X^{(k)}\}$ be independent, jointly Gaussian, self-similar processes with Hurst exponents $H_k$ and comparable scale. Then for fixed weights $a_k$ independent of state or time, $\sum_k a_k X^{(k)}_t$ is Gaussian and, at long lag, its covariance is dominated by the term(s) with the largest $H_k$; i.e., $H_{\mathrm{agg}} \approx \max_k H_k$.
\end{proposition}

This is a rigorous consequence of summing power-law kernels, but it is not the result needed: $\max_k H_k$ can equal but cannot \emph{exceed} every individual $H_k$, whereas the empirical pattern requires $H_{\text{index}}$ to exceed every single name's $H_i$, common factor included at the naive level. Reaching that stronger conclusion requires positing that the common factor $F$ is smoother than every idiosyncratic $X_i$ \emph{a priori} -- which is precisely \eqref{eq:factor}'s empirical input, not something the aggregation math derives.

\begin{caveat}
This step is circular if presented as a derivation of "index smoother than every name" from first principles: the conclusion is smuggled in via the assumed relative smoothness of $F$, not produced by the aggregation mechanics themselves.
\end{caveat}

\section{The gap, closed for the consequence but not the cause}

\citet{zarhali2025nested} supply the missing piece for the aggregation \emph{consequence}, not the smooth/rough split itself. Writing the index as a weighted average with weights $w_i$, the idiosyncratic contribution to index variance scales as $\sum_i w_i^4\sigma_i^4 \sim N_s^{-3}$ for an equal-weighted basket of $N_s$ names, while the common-factor contribution does not shrink with $N_s$. This yields an explicit, checkable threshold:
\begin{equation}
\beta^4 N_s \gg 10 \label{eq:threshold}
\end{equation}
above which the shared smooth factor dominates and the index's apparent Hurst exponent converges to $H_F$, exceeding every individual $H_i$. This is a genuine cross-sectional-diversification derivation, not curve-fitting of the sign of the effect -- it directly closes the gap our Proposition~1 attempt could not.

\begin{finding}
What remains open is not the aggregation \emph{consequence} (Equation~\eqref{eq:threshold}, now well-supported) but the aggregation \emph{cause}: why should the common/market factor be smoother than idiosyncratic volatility in the first place? \citet{zarhali2025nested} import this from \citet{wu2022rough} as data, not as a derived quantity.
\end{finding}

\subsection{A ruled-out shortcut: CLT-of-CLTs}

One might hope the smooth/rough split is itself a byproduct of the same temporal aggregation mechanism (\`a la \citealp{eleuch2018microstructural}) that already produces $H>0$ from microstructure -- i.e., that ``more aggregation'' at the cross-sectional level simply extends ``more aggregation'' at the temporal level.

\begin{proposition}[Cross-sectional aggregation is selection, not smoothing]
Let $\omega_i(t) = \beta_i F(t) + X_i(t)$ with $X_i$ i.i.d.\ across $i$, independent of $F$, each with the same Hurst exponent $H_X$. Then $\frac1N\sum_i \omega_i(t) \to \bar\beta F(t)$ in mean-square as $N\to\infty$, and the roughness exponent of the limit is exactly $H_F$ -- unaffected by $H_X$ and unaffected by $N$ once $N$ is large enough for the idiosyncratic term to vanish. Cross-sectional averaging does not alter or interpolate any component's own Hurst exponent; it only reweights their relative amplitudes.
\end{proposition}

\begin{finding}
\labelnegative. A ``CLT-of-CLTs'' argument fails: cross-sectional aggregation is a pure variance-reduction/selection mechanism. It reveals whichever component already has the largest surviving amplitude; it cannot manufacture smoothness that is not already present in $F$. Moreover, the temporal mechanism that already produces $H>0$ (nearly-critical Hawkes processes, \citealp{eleuch2018microstructural}) does not get smoother with more aggregated order flow -- the roughness exponent there is set by the kernel's tail index, not by sample size. So the cause of $F$'s smoothness cannot be reduced to ``more of the same'' temporal-CLT mechanism; it requires a genuinely different explanation.
\end{finding}

\subsection{A speculative microstructural alternative}

Combining the Mixture-of-Distributions Hypothesis \citep{clark1973subordinated}, the Hawkes-to-rough-volatility mapping \citep{eleuch2018microstructural,jaisson2016rough}, and Grigelionis-type point-process superposition theory \citep{grigelionis1963}, we advance the following hypothesis, original to this investigation and \emph{not} found stated elsewhere in the literature searched:

\begin{conjecture}
Single-name order flow is generated by a comparatively small, tightly-coupled set of participants, pushing its effective Hawkes branching ratio close to criticality and hence toward small $H$ via the \citet{eleuch2018microstructural} mechanism. Market-wide flow superposes many weakly-correlated local feedback loops plus genuinely diffusive macro information, diluting effective endogeneity at the aggregate level and yielding a larger effective $H_F$.
\end{conjecture}

This produces concrete, falsifiable predictions: single-name Hawkes branching ratios should exceed the index's at matched time resolution, and the common/idiosyncratic $H$-gap should track algorithmic-trading penetration. We flag it as genuinely uncertain: \citet{filimonov2015quantifying} estimate E-mini S\&P~500 order flow itself to be highly endogenous ($>70\%$), which directly complicates the naive version of this hypothesis.

\confidence{\labelspeculative{} (Conjecture~1); untested against the cited counter-evidence.}

\section{Evaluation}

\subsection{Simulation of the aggregation mechanism}

We independently simulated \eqref{eq:factor} with $H_F=0.15$, $\beta=0.6$, and i.i.d.\ log-covariance idiosyncratic components ($H_X\approx0.02$), via Cholesky factorization of the respective covariance matrices, and estimated the aggregate's Hurst exponent by structure-function slope as $N_s$ grows.

\begin{table}[h]
\centering
\begin{tabular}{rrr}
\toprule
$N_s$ & $\beta^4 N_s$ & Estimated aggregate $H$ \\
\midrule
1 & 0.13 & 0.041 \\
10 & 1.30 & 0.095 \\
30 & 3.89 & 0.121 \\
100 & 12.96 & 0.140 \\
1000 & 129.6 & 0.148 \\
\bottomrule
\end{tabular}
\caption{Simulated aggregate roughness converges toward the common factor ($H_F=0.15$) as the diversification threshold \eqref{eq:threshold} is crossed.}
\end{table}

\begin{finding}
The mechanism holds up under independent simulation from scratch, not merely as an internal derivation within \citet{zarhali2025nested}: the estimated aggregate $H$ tracks the $\beta^4N_s\gg10$ threshold closely, already at $0.095$ by $N_s=10$ (where $\beta^4N_s=1.3$) and essentially converged by $N_s=100$.
\end{finding}

\subsection{Real-data check}

We estimated Hurst exponents from fifteen years of daily OHLC data (S\&P 500 index and 25 large-cap single names) using a Parkinson range-based log-volatility proxy and a structure-function-slope estimator. The index's $\hat H=0.094$ exceeded every single name (range $0.047$--$0.076$), replicating the qualitative pattern on data none of the cited papers used. As documented in the companion paper on distributional versus temporal multifractality (this series, Paper~1), a more careful nested boundary test applied to the same data complicates this: most series, including the index, are statistically close to the $H=0$ boundary once penalized for the extra parameter, so the absolute magnitudes here should be read cautiously even as the relative ordering (index $>$ every single name) held up under the simple estimator.

\section{Findings summary}

\begin{itemize}
\item \labelestablished{} that classical aggregation theory permits altered memory/self-similarity under heterogeneous mixing (\citealp{granger1980long}; general mechanism family).
\item \labelplausible{}, now closed to \labelestablished{}-adjacent status conditional on its input: the specific diversification threshold \eqref{eq:threshold} of \citet{zarhali2025nested}, independently re-simulated and confirmed here.
\item \labelnegative{} for a naive first-principles derivation of ``index smoother than every name'' from linear mixing alone (Proposition~1's circularity).
\item \labelnegative{} for the CLT-of-CLTs shortcut to explaining common-factor smoothness (Proposition~2).
\item \labelspeculative{} for the proposed microstructural cause (Conjecture~1), with a specific documented counter-example in the literature.
\end{itemize}

\section{Future directions}

\begin{enumerate}
\item Test the branching-ratio prediction of Conjecture~1 directly: compare estimated Hawkes near-criticality for single names versus the index at matched time resolution, if tick-level data with participant attribution becomes available.
\item Re-run the real-data check of Section~5.2 with true high-frequency realized variance rather than the daily Parkinson proxy, to see whether the aggregation pattern survives the nested boundary test once measurement noise is reduced.
\item Investigate whether the index's effective decorrelation timescale $T$ also grows with $N_s$ under \eqref{eq:threshold}, as suggested by a cross-pollination result in the companion SPX/VIX paper (this series, Paper~4) -- this would connect the aggregation mechanism directly to a pricing-relevant consequence.
\item Attempt a genuine (non-circular) first-principles derivation of $H_F>H_i$ starting from a specific, falsifiable economic microfoundation, rather than importing the split as data.
\end{enumerate}

\section*{Acknowledgments}
This paper synthesizes research conducted by an orchestrated multi-agent AI investigation, with all citations independently re-verified against primary sources and all numerical claims re-run and confirmed before inclusion here.

\bibliographystyle{plainnat}
\begin{thebibliography}{99}

\bibitem[Bennedsen et al.(2022)]{bennedsen2022decoupling}
M.~Bennedsen, A.~Lunde, and M.~S.~Pakkanen.
\newblock Decoupling the short- and long-term behavior of stochastic volatility.
\newblock \emph{Journal of Financial Econometrics}, 20(5):961--1006, 2022.

\bibitem[Bolko et al.(2023)]{bolko2023gmm}
A.~E.~Bolko, K.~Christensen, M.~S.~Pakkanen, and B.~Veliyev.
\newblock A GMM approach to estimate the roughness of stochastic volatility.
\newblock \emph{Journal of Econometrics}, 235(2):745--778, 2023.

\bibitem[Clark(1973)]{clark1973subordinated}
P.~K.~Clark.
\newblock A subordinated stochastic process model with finite variance for speculative prices.
\newblock \emph{Econometrica}, 41(1):135--155, 1973.

\bibitem[El~Euch et al.(2018)]{eleuch2018microstructural}
O.~El~Euch, M.~Fukasawa, and M.~Rosenbaum.
\newblock The microstructural foundations of leverage effect and rough volatility.
\newblock \emph{Finance and Stochastics}, 22(2):241--280, 2018.

\bibitem[Filimonov and Sornette(2015)]{filimonov2015quantifying}
V.~Filimonov and D.~Sornette.
\newblock Quantifying reflexivity in financial markets: toward a prediction of flash crashes.
\newblock \emph{Physical Review E}, 85(5):056108, 2015.

\bibitem[Gatheral et al.(2018)]{gatheral2018volatility}
J.~Gatheral, T.~Jaisson, and M.~Rosenbaum.
\newblock Volatility is rough.
\newblock \emph{Quantitative Finance}, 18(6):933--949, 2018.

\bibitem[Granger(1980)]{granger1980long}
C.~W.~J.~Granger.
\newblock Long memory relationships and the aggregation of dynamic models.
\newblock \emph{Journal of Econometrics}, 14(2):227--238, 1980.

\bibitem[Grigelionis(1963)]{grigelionis1963}
B.~Grigelionis.
\newblock On the convergence of sums of random step processes to a Poisson process.
\newblock \emph{Theory of Probability \& Its Applications}, 8(2):177--182, 1963.

\bibitem[Jaisson and Rosenbaum(2016)]{jaisson2016rough}
T.~Jaisson and M.~Rosenbaum.
\newblock Rough fractional diffusions as scaling limits of nearly unstable heavy tailed Hawkes processes.
\newblock \emph{Annals of Applied Probability}, 26(5):2860--2882, 2016.

\bibitem[Wu et al.(2022)]{wu2022rough}
Q.~Wu, J.-F.~Muzy, and E.~Bacry.
\newblock From rough to multifractal volatility: the log S-fBM model.
\newblock \emph{Physica A}, 604:127919, 2022. arXiv:2201.09516.

\bibitem[Zaffaroni(2004)]{zaffaroni2004aggregation}
P.~Zaffaroni.
\newblock Contemporaneous aggregation of linear dynamic models in large economies.
\newblock \emph{Journal of Econometrics}, 120(1):75--102, 2004.

\bibitem[Zarhali et al.(2025)]{zarhali2025nested}
O.~Zarhali, C.~Aubrun, E.~Bacry, J.-P.~Bouchaud, and J.-F.~Muzy.
\newblock A nested factor model for equity markets: reconciling multifractal stock returns and rough index volatilities.
\newblock \emph{arXiv:2505.02678}, 2025.

\end{thebibliography}

\end{document}
